\section*{Chapter \ref{ch:fm:the-risk-management-function} --- The Risk-Management Function}

\begin{solution}{exo:fm:the-risk-management-function:1}
Severity 1 (the head of desk, the same day); severity 2 (the \omterm{def:fm:the-risk-management-function:cro}{chief risk officer}, within a business day); severity 3 (the \omterm{def:fm:the-risk-management-function:cro}{risk committee}, within
two business days).
\end{solution}

\begin{solution}{exo:fm:the-risk-management-function:2}
From \$2.3 to \$4.0 billion: 74\%. The 2008 limit was \$1.5 billion, or 60\%, above the \$2.5 billion of the previous year's method.
\end{solution}

\begin{solution}{exo:fm:the-risk-management-function:3}
$(5+1)/2=3$ business days; $(10+1)/2=5.5$ with meetings every tenth day.
\end{solution}

\begin{solution}{exo:fm:the-risk-management-function:4}
It rewrites the record: the breaches between the effective day and the approval day disappear from a register that uses the effective limit, so
no one can later see that the firm was over its limit or that the increase was granted to cover a breach. A register must keep limits as known
each day, with approval and effective dates.
\end{solution}

\begin{solution}{exo:fm:the-risk-management-function:5}
Desks that ask for an increase wait for the next meeting, three business days on average, with the breach open. When they must cut while they
wait, most breaches close before the meeting and the request lapses: the mean duration falls from 1.64 to 1.25 days.
\end{solution}

\begin{solution}{exo:fm:the-risk-management-function:6}
Two people are involved, but the approver is not independent: the head of desk is paid on the desk's results and wants the same outcome as the
trader. The second pair of eyes must belong to someone whose interests differ, such as the risk function or the committee.
\end{solution}

\begin{solution}{exo:fm:the-risk-management-function:7}
Five breaches, at $-51$, $-27$, $-9$, $-5$ and $-2$ business days from the change. The one at $-27$ is flagged for a limit raised during a
breach; the one at $-5$, when the first increase expired, for a limit raised and repeated increases; the one at $-2$, closed by the model change, for
a model changed during a breach and repeated increases. The other two closed by cuts and are not flagged.
\end{solution}

\begin{solution}{exo:fm:the-risk-management-function:8}
The head of trading is the head of the businesses whose risks are measured: the objection runs through the person it objects to, whose pay
depends on the revenue the limits constrain. Independence needs a reporting line to the \omterm{def:fm:the-risk-management-function:cro}{risk committee} and the chief executive, not expertise in
the line manager.
\end{solution}

\begin{solution}{pb:fm:the-risk-management-function:1}
\begin{enumerate}
\item See \cref{def:fm:the-risk-management-function:cro}.
\item A reporting line outside the businesses, pay independent of their revenue, and direct access to the \omterm{def:fm:the-risk-management-function:cro}{risk committee}.
\item Independence and authority combined with direct interaction with business managers, early escalation (some from the summer of 2006), and
information not filtered by hierarchy.
\item See \cref{def:fm:the-risk-management-function:foureyes}; limit increases and changes to risk models.
\item \$2.3 billion; \$3.3 billion for fiscal 2007 (end of 2006); \$3.5 billion on 7 September 2007; \$4.0 billion approved on 14 January 2008,
effective 3 December 2007.
\item It removed from the record the excesses over the \$3.5 billion limit between 3 December and the approval.
\item Commercial real estate, private equity and, for a time, leveraged loans; a \$2.3 billion bridge equity position, from May to August 2007.
\item Insufficient evidence that conduct was outside the business judgement rule or reckless: unwise decisions were still management's to make.
The risk function's authority must be written into procedure beforehand.
\item See \cref{def:fm:the-risk-management-function:escalation}.
\item Severity 1 up to 5\% and two days, 3 above 20\% or five days, 2 between; for example 3\%/1 day, 10\%/1 day and 3\%/6 days give 1, 2 and 3.
\item See \cref{prop:fm:the-risk-management-function:wait}.
\item 45.3\%, 14.4\% and 5.4\%.
\item 1.15, 1.64 and 1.25 days; 31.8, 43.5 and 34.6 breach-days a year; 2.44\%, 1.33\% and 0.74\% of desk-days.
\item Limits raised or models changed during a breach, and repeated increases; 54.7\%, 23.1\% and 11.9\%.
\item Most breaches are one-day drifts that reverse whatever the regime; the mean duration, the breach-days and the severity counts show the
difference.
\item A temporary increase that expires and is renewed at once, then a breach of 7.7\% closed by a model change that lowers the measured risk by a
fifth without any change in the position.
\item Limits and models as known each day, with approval and effective dates for every change, and every flag, never cleared.
\item See \cref{def:fm:the-risk-management-function:culture}; the share of breaches closed by increases, the flagged share, exceptions past their
expiry.
\item Head of desk approving: 45.3\% of breaches closed by raising the limit, median one day (mean 1.15). Committee approving: 14.4\%, median one
day but mean 1.64 and more breach-days; with cutting while waiting, 5.4\% and a mean of 1.25.
\item Increases and model changes need the \omterm{def:fm:the-risk-management-function:cro}{risk committee}'s approval, with the desk cutting towards its limit while it waits; every change is
recorded with its approval date, and repeated increases and changes during a breach go to the committee as red flags.
\end{enumerate}
\end{solution}

\begin{solution}{iq:fm:the-risk-management-function:1}
To the \omterm{def:fm:the-risk-management-function:cro}{risk committee} of the board and to the chief executive, not to the head of the businesses whose risk it measures, so that objections cannot
be filtered or overruled by those they concern.
\iqlookfor{independence of the reporting line and direct access to the board.}
\end{solution}

\begin{solution}{iq:fm:the-risk-management-function:2}
It is logged and graded; the head of desk is told the same day; the desk cuts or asks for a temporary increase approved independently, with an
expiry; if it persists, it is escalated to the \omterm{def:fm:the-risk-management-function:cro}{chief risk officer}.
\iqlookfor{the procedure, not only the cut.}
\end{solution}

\begin{solution}{iq:fm:the-risk-management-function:3}
Refuse to renew by default and take it to the committee: either the limit is wrong and should be reset through the annual process, or the desk is
using increases to carry more risk than it was given.
\iqlookfor{repeated increases as a red flag.}
\end{solution}

\begin{solution}{iq:fm:the-risk-management-function:4}
The breach stands under the old model until the new one is independently validated; validation checks the model on its merits, not on its effect,
and the change is recorded with its date and flagged.
\iqlookfor{validation independent of the breach.}
\end{solution}

\begin{solution}{iq:fm:the-risk-management-function:5}
The share of breaches closed by increases, changes made during breaches, backdated changes, exceptions past expiry, and the time for bad news to
reach the committee.
\iqlookfor{patterns, not counts.}
\end{solution}

\begin{solution}{iq:fm:the-risk-management-function:6}
Append-only tables of limits, model versions and exposures, each with effective and recorded timestamps; breaches derived from them as of each
day; approvals as separate signed records; flags stored and never deleted.
\iqlookfor{bitemporal records and append-only storage.}
\end{solution}
